\section{Bounded paths exclude AUC renormings}\label{rec:obstruction}

The midpoint gain controls two opposite perturbations together.
To distinguish it from one-sided asymptotic convexity, we use positive
path vectors that remain bounded as their paths grow. At each stage,
the next coordinate can be chosen from a weakly null sequence of
increments whose norms stay bounded away from zero.

\begin{lemma}[Bounded trees]\label{rec:bounded-tree}
Let $X$ be a Banach space. Suppose that for every $h\ge1$ there are
vectors $u_s^{(h)}$, indexed by $s\in\N^{\le h}$, such that
$0<m\le\|u_s^{(h)}\|\le M<\infty$ uniformly. Suppose also that for
$|s|<h$ the increments
$d_{s,k}^{(h)}=u_{s^\frown k}^{(h)}-u_s^{(h)}$ are weakly null as
$k\to\infty$ and have norms at least $c>0$, uniformly in $h,s,k$.
If $N$ is a norm and $0<\alpha\le\beta<\infty$ satisfy
$\alpha\|\cdot\|\le N\le\beta\|\cdot\|$, then
\[
 \overline\delta_N\left(\frac{\alpha c}{2\beta M}\right)=0.
\]
The conclusion also holds if one infinite tree has these properties.
\end{lemma}

\begin{proof}
Set $t_0=\alpha c/(2\beta M)$ and suppose
$0<\gamma<\overline\delta_N(t_0)$. At a current vector $u_s^{(h)}=u$,
choose a closed finite-codimensional subspace $F$ such that
\[
 N(u+t_0N(u)y)\ge(1+\gamma)N(u)
 \qquad(y\in F,\ N(y)=1).
\]
Convexity extends this inequality to every $w\in F$ with
$N(w)\ge t_0N(u)$. More explicitly, for a unit center $x$ and a unit
direction $y$, convexity at the radii $0,t_0,r$ gives
$N(x+ry)-1\ge(r/t_0)(N(x+t_0y)-1)$ for $r\ge t_0$.

Weak nullity persists under an equivalent norm. In the finite-dimensional
quotient $X/F$ the increments converge to zero in norm, so there are
$w_k\in F$ with $N(w_k-d_{s,k}^{(h)})\to0$.
For large $k$,
$N(w_k)\ge\alpha c/2\ge t_0N(u)$.
Choose $k$ with error less than $\gamma N(u)/2$. The preceding estimate
then gives
\[
 N(u_{s^\frown k}^{(h)})\ge(1+\gamma/2)N(u_s^{(h)}).
\]
Choose $h$ such that $\alpha m(1+\gamma/2)^h>\beta M$ and iterate this
choice for $h$ steps, a contradiction. Restricting an infinite tree to
that height gives the same argument. Finally,
$\overline\delta_N(t_0)\ge0$: intersect a finite-codimensional subspace
with the kernel of a functional norming a given center. Both signs of
every perturbation in that kernel have norm at least that of the center.
\end{proof}

\begin{proof}[Proof of the renorming assertion in Theorem~\ref{rec:main}]
On the infinite tree define
$v_s=\sum_{r\preceq s}e_r$. Along this path the fields of $v_s$ are
$P=1,Q=0$, so $N_\Sigma(v_s)=1$. Its child increments are $e_{sk}$,
and for a finite scalar family $(a_k)$ the recursion gives
\begin{equation}\label{rec:sibling-formula}
 N_\Sigma\Bigl(\sum_k a_ke_{sk}\Bigr)
 =2\max\{\|(a_k^+)\|_2,\|(a_k^-)\|_2\}
 \le2\|(a_k)\|_2.
\end{equation}
Here $a^+=\max\{a,0\}$ and $a^-=\max\{-a,0\}$.
In particular $N_\Sigma(e_{sk})=2$.
The map taking the unit vectors of $\ell_2$ to $e_{sk}$ extends to a
bounded linear map by \eqref{rec:sibling-formula}. Since the unit vectors
of $\ell_2$ are weakly null, so are these increments against every
continuous linear functional on $X_\Sigma$.
Lemma~\ref{rec:bounded-tree} applies with $m=M=1,c=2$.

For a node $s\in T_n$, let $v_s$ be the indicator of the path from
$r_n$ to $s$, with joining coordinate zero. Its component fields are
again $P=1,Q=0$ along the path. At $\rho$ they both equal one. Thus
\[
 N_0(v_s)=N_0(e_s)=1,\qquad N_J(v_s)=N_J(e_s)=2.
\]
At a fixed non-leaf $s$ in a component, the sibling sequence is weakly
null: restricting any functional to that component gives a continuous
functional on a space equivalent to $\ell_2(T_n)$ by
\eqref{rec:block-equivalence}. Alternatively, the sibling computation
in \eqref{rec:sibling-formula} gives a direct bounded map from $\ell_2$;
for $N_0$ its factor two is replaced by one.
The joined trees contain components of every finite height.
Lemma~\ref{rec:bounded-tree} therefore applies to $X_0$ with
$m=M=c=1$, and to $X_J$ with $m=M=c=2$.
In both cases the vanishing one-sided modulus occurs at
$t_0=\alpha/(2\beta)$; for $X_\Sigma$ it occurs at
$t_0=\alpha/\beta$.
\end{proof}

The argument permits the finite-codimensional subspace to change at
every node. Its contradiction uses only a fixed positive gain and
arbitrarily long bounded paths. Reflexivity of the joined constructions
is compatible with this obstruction; it does not imply the existence
of an AUC renorming.
